\section{Marco referencial}
\label{sec:marco_referencial}

Dos tipos de referentes delimitan esta investigación: el marco curricular colombiano vigente para la educación media, y los referentes internacionales en diseño de tareas matemáticas que orientan sus decisiones pedagógicas.

\subsection{Marco curricular colombiano}

Como objeto matemático, la trigonometría no se agota en el triángulo rectángulo. También articula medición angular, periodicidad y modelación de fenómenos cíclicos, por lo que conecta la geometría con el análisis matemático y con múltiples aplicaciones científicas y tecnológicas.

Los \emph{Lineamientos Curriculares de Matemáticas} ubican la trigonometría en relación con el pensamiento variacional, la modelación y el análisis de fenómenos de cambio \parencite{ministeriodeeducacionnacionalLineamientosCurriculares1998}. Los \emph{Estándares Básicos de Competencias} van más allá: en décimo y undécimo, los estudiantes deben interpretar el comportamiento de funciones, relacionar parámetros algebraicos con representaciones gráficas y modelar situaciones periódicas del mundo real \parencite{ministeriodeeducacionnacionalEstandaresBasicos2006}. Y los \emph{Derechos Básicos de Aprendizaje} cierran el círculo insistiendo en la comprensión funcional y en el uso flexible de representaciones, no solo en la ejecución de procedimientos algorítmicos \parencite{ministeriodeeducacionnacionalDerechosBasicos2016}.

Ese marco normativo sustenta la decisión de centrar el estudio en el pensamiento trigonométrico y en el diseño de tareas con integración de IA generativa. No es solo normativo, sin embargo. También incorpora estudios que documentan dificultades recurrentes en trigonometría escolar —el tránsito de razón trigonométrica a función, la interpretación de la periodicidad, la articulación de registros semióticos \parencite{delgadoEstrategiaDidacticaPara2023, chaconComprensionRazonesTrigonometricas2007, zubietaTipificacionErroresDificultades2018}—, y esos hallazgos operan como criterio para seleccionar y diseñar las tareas de la investigación.

\subsection{Referentes internacionales en diseño de tareas matemáticas}

¿Qué hace que una tarea matemática genere aprendizaje real, en vez de solo ocupar el tiempo de la clase? Esa pregunta atraviesa la literatura internacional sobre diseño de tareas, y tres referentes la responden desde ángulos distintos —cada uno completando lo que el anterior deja sin resolver.

\textbf{NCTM (2014)} ofrece una primera respuesta, centrada en la demanda cognitiva. \emph{Principles to Actions} \parencite{nctmPrinciplesActions2014} distingue cuatro niveles —memorización, procedimientos sin conexión, procedimientos con conexión, y hacer matemáticas— y sostiene que la instrucción de calidad exige sostenerse en los dos niveles superiores. Es un criterio útil. Pero deja una pregunta sin responder: ¿de dónde sale ese espacio de exploración que una tarea de alta demanda necesita para funcionar? Este referente sustenta, de todos modos, la distinción ejercicio/tarea que articula el marco conceptual de la investigación.

Ahí entra \textbf{Watson y Ohtani (2015)}. El estudio de la ICMI sobre diseño de tareas \parencite{watsonTaskDesignMathematics2015} sintetiza décadas de investigación y precisa algo que el criterio de NCTM da por sentado: ese espacio de posibilidades matemáticas no está en el enunciado de la tarea. Está en las herramientas disponibles, el tiempo, las interacciones con pares y con el docente, los recursos tecnológicos. Para esta investigación, esa precisión es directamente operativa: ChatGPT y GeoGebra no son accesorios del enunciado. Son parte de ese espacio.

Falta todavía una pieza: ¿qué es exactamente lo que ese espacio de exploración intenta mover en la cabeza del estudiante? \textbf{Sfard (1991)} \parencite{sfardDualNatureMathematical1991} responde con la distinción entre concepción operacional —el objeto matemático como proceso o algoritmo— y estructural —como objeto estático, con propiedades. El tránsito entre ambas es uno de los desafíos cognitivos más importantes del aprendizaje matemático avanzado, y la trigonometría lo ilustra con particular claridad: el seno puede ser un procedimiento para calcular razones en un triángulo, o puede ser una función —periódica, continua, diferenciable— con dominio en los reales. Las tareas de esta investigación se diseñan para mover al estudiante de lo primero hacia lo segundo.

Los tres referentes no compiten entre sí: se completan. NCTM fija el nivel de exigencia. Watson y Ohtani ubican dónde vive ese espacio de exigencia. Sfard nombra la transformación cognitiva que se busca dentro de él.
